\chapter{Conclusiones y Trabajos Futuros}

\section{Conclusiones}

El desarrollo de la presente tesis permite concluir que la solución propuesta constituye una alternativa tecnológicamente coherente y académicamente defendible para mejorar la mediación cultural en museos peruanos mediante una arquitectura BYOD, reconocimiento contextual por proximidad BLE y una capa conversacional sustentada por recuperación aumentada. A diferencia de una audioguía convencional, la solución planteada no se limita a reproducir contenido estático, sino que integra contexto físico, interacción multimodal y conocimiento curatorial estructurado para producir respuestas más situadas durante la visita.

En términos de problema y oportunidad, la investigación confirma que existe un espacio real para una solución de esta naturaleza dentro del segmento de museos estatales de alta afluencia. El análisis de mercado permitió delimitar un mercado objetivo institucional plausible, mientras que la revisión de la oferta tecnológica mundial mostró que la diferenciación de la solución propuesta no reside en replicar aplicaciones de guía existentes, sino en combinar proximidad indoor, uso del smartphone del visitante, voz e inteligencia contextual dentro de una misma experiencia.

Desde la perspectiva técnica, la tesis demuestra que el núcleo funcional del sistema puede materializarse con una arquitectura distribuida y de bajo costo. La combinación de beacons BLE sobre ESP32, aplicación móvil y backend RAG permite sostener una lógica de mediación contextual compatible con el entorno museístico. En particular, el trabajo realizado confirma que la granularidad por sala o nodo cercano es un objetivo realista y suficiente para activar contenido pertinente sin recurrir a infraestructuras de localización de alta precisión, más costosas y complejas de desplegar.

También puede concluirse que la estructuración explícita de los datos de entrada es una condición central para el éxito del proyecto. La tesis dejó establecido que la utilidad de la solución propuesta depende de articular, de manera trazable, identificadores BLE, señales del smartphone, eventos de interacción y corpus curatoriales con metadatos consistentes. Esta decisión metodológica evita tratar a la capa conversacional como un chatbot genérico y refuerza su naturaleza de sistema contextualizado y anclado en evidencia institucional.

Otra conclusión importante es que el prototipo ya desarrollado supera el nivel de formulación conceptual. Los repositorios técnicos del sistema demuestran que existe una implementación operativa distribuida que cubre emisión BLE, experiencia móvil, persistencia local y respuesta contextual. En ese sentido, la tesis no se limita a proponer una idea: documenta una base real de experimentación, integración y validación que conecta la arquitectura con un prototipo reproducible.

Desde la perspectiva económica y financiera, la propuesta también resulta defendible. El modelo construido a partir de paquetes institucionales Básico, Estándar y Avanzado muestra que la solución propuesta puede organizarse como un servicio híbrido, con ingreso inicial por implementación y flujos recurrentes por mantenimiento, IA y analítica. Bajo el escenario base de evaluación a diez años, con 35 museos captados progresivamente, el proyecto presenta un VAN positivo de S/444\,585, una TIR de 33.99\%, un índice de rentabilidad de 2.48 y una recuperación descontada de 6.35 años, lo que indica viabilidad económica bajo supuestos prudentes. La conclusión no es que el proyecto esté exento de riesgo, sino que su sostenibilidad mejora cuando la solución se transforma en una cartera recurrente, cuando se protege la cobranza institucional y cuando el costo de la capa conversacional se gestiona con disciplina contractual.

Asimismo, la tesis permite concluir que la viabilidad de la solución propuesta no depende solo del desempeño técnico o del precio de implementación, sino de la capacidad de articular estándares de calidad, seguridad e interoperabilidad desde etapas tempranas. La incorporación de referencias como ISO/IEC/IEEE 12207, Bluetooth Core Specification, RFC 8259, OWASP, ISO/IEC 25012, CIDOC CRM, LIDO, Dublin Core y la Ley N.\ 29733 refuerza que el proyecto ya se ha pensado bajo un marco de gobernanza técnica compatible con una futura evolución institucional.

Finalmente, el trabajo también deja explícitos sus límites. La validación desarrollada corresponde todavía a un prototipo de tesis y no a un despliegue museal completo. La cobertura BLE se ha trabajado como prueba de sala y nodo; la aplicación prioriza Android en entorno de \textit{Development Build}; el corpus RAG sigue acotado a una maqueta curatorial específica; y el backend aún opera en red local con herramientas adecuadas para experimentación, pero no equivalentes a un servicio de producción. Sin embargo, estas limitaciones no invalidan la propuesta. Más bien, delimitan con claridad el punto alcanzado y refuerzan el valor del proyecto como base sólida para escalamiento y transferencia tecnológica.

\section{Relación entre objetivos y resultados}

La Tabla~\ref{tab:objetivos-resultados-conclusiones} resume la relación entre los objetivos desarrollados a lo largo de la tesis y los resultados alcanzados. Esta lectura permite cerrar el documento de manera trazable, conectando el problema inicial, el diseño técnico, la validación del prototipo y la evaluación económica.

\begin{table}[!ht]
    \centering
    \scriptsize
    \setlength{\tabcolsep}{4pt}
    \resizebox{\textwidth}{!}{%
    \begin{tabular}{>{\raggedright\arraybackslash}p{4.1cm} >{\raggedright\arraybackslash}p{6.0cm} >{\raggedright\arraybackslash}p{5.8cm}}
        \hline
        \textbf{Objetivo o necesidad} & \textbf{Resultado alcanzado} & \textbf{Evidencia principal} \\
        \hline
        Identificar el problema de mediación digital en museos peruanos. & Se caracterizó la brecha entre experiencia presencial, acceso a contenido contextual y limitaciones de mediación tradicional. & Capítulos 1 y 2: problema, panorama nacional, demanda y estado del arte. \\
        Delimitar un mercado objetivo institucional. & Se definió un segmento de museos estatales y una cartera comercial por paquetes. & Capítulos 3, 4 y 5: mercado objetivo, oferta tecnológica e ingresos proyectados. \\
        Diseñar una solución técnica viable y modular. & Se formuló una arquitectura con BLE, app móvil, persistencia local, backend RAG y corpus curatorial. & Capítulos 7, 8 y 9: tecnologías, datos de entrada, arquitectura e interfaces. \\
        Validar el núcleo funcional del prototipo. & Se documentaron pruebas de reconocimiento de contexto, consulta contextual, fallback y despliegue reproducible. & Capítulo 9: protocolo de validación, pseudocódigo y trazabilidad metodológica. \\
        Evaluar sostenibilidad económica. & El escenario base obtuvo VAN positivo, TIR superior a la tasa de corte e índice de rentabilidad mayor que uno. & Capítulo 10: flujo financiero, sensibilidad e indicadores económicos. \\
        Alinear el diseño con normas y buenas prácticas. & Se identificaron estándares aplicables y controles mínimos para una futura transferencia institucional. & Capítulo 11: normas técnicas, matriz de verificación y alcance real de cumplimiento. \\
        \hline
    \end{tabular}%
    }
    \caption[Relación entre objetivos y resultados]{Relación entre objetivos de la tesis, resultados alcanzados y evidencias principales. Fuente: elaboración propia.}
    \label{tab:objetivos-resultados-conclusiones}
\end{table}

Como se observa, la tesis no concluye únicamente con una propuesta conceptual. El resultado integrado combina análisis de mercado, diseño técnico, evidencia de prototipo, evaluación financiera y criterios normativos. Esta articulación es la que permite sostener la pertinencia de la solución como proyecto de ingeniería aplicada.

\section{Trabajos futuros}

La primera línea de trabajo futuro consiste en ampliar la validación del sistema en condiciones museales más cercanas a la operación real. Esto implica desplegar la solución propuesta en un conjunto mayor de salas, aumentar el número de beacons instalados, calibrar la infraestructura en recorridos completos y contrastar su comportamiento con mayor diversidad de obstáculos físicos, flujo de visitantes y condiciones de uso.

Una segunda línea prioritaria es profundizar la evaluación con usuarios. La tesis ha validado coherencia técnica y funcional del prototipo, pero una etapa posterior debería incorporar pruebas controladas con visitantes y personal del museo para medir comprensión, utilidad percibida, accesibilidad, carga cognitiva, aceptación de la interacción por voz y efecto de la contextualización sobre la experiencia de visita.

Un tercer frente de trabajo futuro es fortalecer la capa conversacional y el corpus curatorial. Resulta necesario ampliar la cobertura de salas, obras y documentos, mejorar la segmentación e indexación del contenido, enriquecer los metadatos curatoriales y comparar el desempeño de distintos modelos y estrategias de recuperación. También sería pertinente incorporar mecanismos más explícitos de evaluación de respuestas, cobertura de fuentes y control de alucinación en contextos patrimoniales.

La cuarta línea se relaciona con la evolución hacia una arquitectura institucional más robusta. En fases posteriores, la solución propuesta debería migrar desde un backend local de laboratorio hacia un despliegue con mayor disponibilidad, observabilidad, endurecimiento de seguridad y gestión operativa. Esto incluye autenticación más formal entre servicios, cifrado en tránsito fuera de entornos controlados, políticas de retención y consentimiento para datos de voz, y documentación de API apta para operación sostenida.

La quinta línea futura consiste en ampliar la interoperabilidad patrimonial del sistema. Aunque el prototipo ya adopta principios de metadatos mínimos consistentes, una evolución posterior podría profundizar la alineación con esquemas como CIDOC CRM, LIDO o Dublin Core para facilitar integración con catálogos museales más formales, interoperabilidad entre instituciones y reutilización del contenido en otros entornos culturales.

Finalmente, una línea de trabajo especialmente valiosa sería explorar la expansión comercial y territorial del modelo. La solución propuesta podría extenderse a otros museos del mismo clúster regional, incorporar nuevas rutas, exposiciones temporales o públicos específicos, e incluso ensayar variantes de monetización complementaria. Esa expansión debería hacerse manteniendo el principio que atraviesa toda la tesis: crecer desde pilotos controlados, con evidencia técnica, económica y operativa suficiente antes de escalar.

\section{Recomendaciones}

Para una siguiente etapa, se recomienda iniciar con un piloto institucional acotado a pocas salas, pero con protocolo formal de pruebas, consentimiento informado para datos de interacción y métricas de experiencia de usuario. Este piloto debería medir no solo si la tecnología funciona, sino si mejora la comprensión de las piezas, reduce fricción durante la visita y aporta información útil al personal del museo.

También se recomienda separar claramente el entorno de laboratorio del entorno de producción. El backend local y los modelos usados en el prototipo son adecuados para investigación y demostración, pero una operación institucional debería incorporar autenticación, cifrado, monitoreo, respaldo del corpus, control de costos de IA y límites explícitos de uso por contrato.

Finalmente, se recomienda mantener la arquitectura modular planteada en la tesis. La mayor fortaleza de la solución no está en una tecnología aislada, sino en la combinación controlada de proximidad BLE, smartphone del visitante, corpus curatorial, voz, recuperación aumentada y evaluación económica. Preservar esa modularidad facilitará corregir, ampliar o sustituir componentes sin perder la lógica integral del sistema.
